\chapter{Cross-cutting concerns: logging, time, ID}
\label{ch:cross-cutting}

\begin{aside}
Some concerns do not belong to a single architectural layer.
Time, logging, ID generation, randomness --- each is touched
by many components, each is a source of bugs if treated
carelessly. This chapter pulls them together and shows the
patterns \maya{} uses for each.
\end{aside}

\section{Time}

A surprising amount of TUI code touches time:

\begin{itemize}[leftmargin=*]
  \item Animations advance over time.
  \item Timeouts fire after a duration.
  \item Timestamps display ``now''.
  \item Debounce windows compare ``now'' to previous.
  \item Profiling measures elapsed time.
\end{itemize}

C++26 gives several clocks:

\begin{itemize}[leftmargin=*]
  \item \code{std::chrono::system\_clock}: wall-clock time.
    Can jump backwards (NTP, user adjustment).
  \item \code{std::chrono::steady\_clock}: monotonically
    increasing. Cannot jump back. Use for durations.
  \item \code{std::chrono::high\_resolution\_clock}: alias for
    one of the above; system-dependent.
\end{itemize}

\subsection*{When to use which}

\begin{itemize}[leftmargin=*]
  \item Display: \code{system\_clock} (the user wants to see
    actual wall-clock).
  \item Animations, timeouts, debounce, profiling:
    \code{steady\_clock}.
\end{itemize}

Mixing them is a bug source. A debounce that uses
\code{system\_clock} can fire repeatedly or never if the clock
jumps.

\subsection*{Wrapping time for testability}

For test determinism, abstract the clock:

\begin{lstlisting}
class Clock {
public:
    virtual std::chrono::steady_clock::time_point now() const = 0;
};

class RealClock : public Clock {
    std::chrono::steady_clock::time_point now() const override {
        return std::chrono::steady_clock::now();
    }
};

class FakeClock : public Clock {
    std::chrono::steady_clock::time_point current_;
public:
    void advance(std::chrono::milliseconds d) { current_ += d; }
    std::chrono::steady_clock::time_point now() const override {
        return current_;
    }
};
\end{lstlisting}

Inject the clock into modules that care; use \code{RealClock}
in production, \code{FakeClock} in tests.

\section{Logging}

Already covered (Chapter~\ref{ch:observability}). The
cross-cutting point: every layer of your app may want to log,
not just the runtime. Make logging accessible:

\begin{itemize}[leftmargin=*]
  \item Pass a \code{Logger\&} via the context.
  \item Inside widgets, accept an optional \code{Logger*}.
  \item For tests, supply a \code{NullLogger}.
\end{itemize}

\subsection*{Avoiding global loggers}

A global \code{spdlog::get(\"app\")} pattern is convenient but
breaks test isolation. The injection pattern is better; one
logger per app instance, one app per test.

\section{IDs}

Generating unique IDs comes up everywhere:

\begin{itemize}[leftmargin=*]
  \item Focus IDs.
  \item Element keys for reconciliation.
  \item Database row IDs.
  \item Correlation IDs for log entries.
\end{itemize}

\subsection*{Patterns}

\begin{enumerate}[leftmargin=*]
  \item \textbf{Monotonic counter.} Simple, safe within one
    process. Not unique across runs.
  \item \textbf{UUID.} Unique across processes, machines, time.
    16 bytes; \code{libuuid} or \code{stduuid}.
  \item \textbf{Hash of content.} Deterministic; same input
    gives same ID. Useful for cache keys.
  \item \textbf{Timestamp + counter.} Sortable IDs.
    Snowflake-style.
\end{enumerate}

\subsection*{Monotonic counter implementation}

\begin{lstlisting}
class IdGenerator {
    std::atomic<uint64_t> next_{1};
public:
    uint64_t next() { return next_.fetch_add(1, std::memory_order_relaxed); }
};
\end{lstlisting}

Thread-safe via atomic. Zero is reserved for ``invalid'';
counts up from 1.

\section{Randomness}

\code{std::rand} is broken (not really random, not thread-safe).
Use \code{<random>}:

\begin{lstlisting}
thread_local std::mt19937 rng{std::random_device{}()};
\end{lstlisting}

Each thread gets its own seeded engine. Use
\code{std::uniform\_int\_distribution} for ranges.

\subsection*{Reproducibility}

For tests, seed the RNG deterministically:

\begin{lstlisting}
std::mt19937 rng{42};
\end{lstlisting}

Same seed, same sequence. The test fails deterministically when
it fails at all.

\subsection*{Cryptographic randomness}

\code{std::mt19937} is not cryptographically secure. For tokens
or session IDs, use the OS:

\begin{itemize}[leftmargin=*]
  \item Linux: \code{getrandom}.
  \item BSD: \code{arc4random}.
  \item Win32: \code{BCryptGenRandom}.
\end{itemize}

The platform abstraction layer should provide a thin wrapper.

\section{Errors}

Cross-cutting because every layer can produce them. Patterns:

\begin{itemize}[leftmargin=*]
  \item Use \code{Expected<T, E>} for recoverable failures.
  \item Use exceptions for unrecoverable bugs (assertion
    failures).
  \item Log every error at the boundary it crossed; do not
    re-log as it propagates.
  \item Surface to the user via the UI, not via stderr.
\end{itemize}

\section{Configuration access}

Already covered (Chapter~\ref{ch:config}). The cross-cutting
point: any layer might need to read config. The pattern:

\begin{itemize}[leftmargin=*]
  \item The model holds the config struct.
  \item \code{view} reads from \code{model.config.\ldots}.
  \item Widgets that need config take it as a parameter.
\end{itemize}

Do not reach into a global; thread the config through.

\section{Caches}

A cache is cross-cutting: it sits between a producer and a
consumer. Patterns:

\begin{itemize}[leftmargin=*]
  \item \textbf{Memoise} pure functions.
    \code{std::unordered\_map<Key, Result>}.
  \item \textbf{LRU caches} for bounded memory.
  \item \textbf{Signal-driven} caches that invalidate when
    inputs change (Chapter~\ref{ch:reactivity-preview}).
\end{itemize}

The trade-off is always memory vs CPU. Profile before adding a
cache; the framework's fingerprint cache often makes app-level
caches unnecessary.

\section{Locking and synchronisation}

\maya{}'s threading model (Chapter~\ref{ch:runtime}) keeps most
synchronisation invisible. Where you need to share state across
threads:

\begin{itemize}[leftmargin=*]
  \item \code{std::atomic} for single primitives.
  \item \code{std::mutex} + \code{std::lock\_guard} for short
    critical sections.
  \item \code{std::shared\_mutex} for read-heavy workloads.
  \item Lock-free queues for high-throughput producer/consumer.
\end{itemize}

The pattern: take the lock for the minimum scope; never call
out (especially not to user code) while holding a lock.

\section{Async cleanup}

Resources that need explicit cleanup (file descriptors, sockets,
GPU handles) should use RAII:

\begin{lstlisting}
class FdGuard {
    int fd_;
public:
    explicit FdGuard(int fd) : fd_{fd} {}
    ~FdGuard() { if (fd_ >= 0) ::close(fd_); }
    FdGuard(const FdGuard&) = delete;
    FdGuard(FdGuard&& o) noexcept : fd_{std::exchange(o.fd_, -1)} {}
};
\end{lstlisting}

The destructor handles cleanup; move-only semantics prevent
double-close. Use \code{std::unique\_ptr} with a custom deleter
for non-trivial resources.

\section{Validation}

Whenever data crosses a trust boundary (network, file, user
input), validate. Patterns:

\begin{itemize}[leftmargin=*]
  \item Validate at the boundary, not deep in the call stack.
  \item Reject early; return \code{Expected<\ldots>::error}.
  \item Log the rejection (so you can spot patterns).
  \item Never silently truncate or coerce; that hides bugs.
\end{itemize}

\section{Telemetry}

If your app reports anonymous usage stats, the patterns from
Chapter~\ref{ch:security} apply: redact, opt-in by default
in some jurisdictions, transparent about what is sent.

The simplest pattern: log a structured event every time a
meaningful action happens; upload batches in
\code{Cmd::isolated\_task}. Privacy and bandwidth both win.

\section{Memory pressure}

When the system is low on memory:

\begin{itemize}[leftmargin=*]
  \item \maya{}'s frame buffers can be shrunk on demand.
  \item Logs can be trimmed.
  \item Caches can be flushed.
\end{itemize}

The Linux kernel signals via cgroups and OOM thresholds; the
app can register a handler. For most apps, this is overkill;
mention it for the few that need it.

\section{Pitfalls}

\begin{warning}
\textbf{Time arithmetic.} Subtracting two
\code{std::chrono::time\_point}s of different clocks does not
compile (good). Subtracting two \code{system\_clock}
timestamps that span a clock jump gives nonsense (bad). Use
\code{steady\_clock} for durations.
\end{warning}

\begin{warning}
\textbf{Global RNG seeded once.} A single global RNG seeded
at startup gives every call the same sequence relative to
program order. Tests become flaky as soon as call order
changes. Use \code{thread\_local} engines or pass a seed.
\end{warning}

\begin{warning}
\textbf{Caching too eagerly.} A cache that never invalidates
holds stale data forever. Pair every cache with an
invalidation rule.
\end{warning}

\begin{warning}
\textbf{Validating in the wrong place.} Validating deep in a
function means callers have already done partial work that
must be undone. Validate at the entry point.
\end{warning}

\section*{Workshop}

\begin{enumerate}[leftmargin=*]
  \item Inject a fake clock into a debounce helper. Test that
    the debounce fires exactly when the clock crosses the
    threshold.
  \item Generate IDs with a thread-safe counter. Verify
    uniqueness under high concurrency.
  \item Wrap a file descriptor in a \code{FdGuard}. Test that
    closing happens deterministically even on early return.
\end{enumerate}

\begin{recap}
Time, logging, IDs, randomness, errors, config, caches, locks,
cleanup, validation, telemetry, memory pressure --- each is a
cross-cutting concern. None belongs to one layer. The patterns:
inject (do not globalise), validate at the boundary, RAII for
resources, atomic or thread-local for shared state. Test with
fakes; production uses the real thing.
\end{recap}
